\section*{Ethical Considerations}\label{sec:ethics}

Our experiments use jailbreak prompt datasets (HarmThoughts, AdvBench) that contain harmful instructions. We use these solely to study model-internal safety mechanisms and do not generate, amplify, or redistribute harmful content. All harmful prompts originate from previously published, publicly available datasets.

This work characterizes how hidden-state probes detect safety-critical transitions during reasoning, including failure modes such as layer selection bias. A potential dual-use risk is that adversaries could exploit knowledge of probe weaknesses---particularly the sensitivity to layer selection---to craft inputs that evade probe-based safety monitors. We believe the benefits of transparently documenting these failure modes outweigh this risk: defenders cannot build robust monitoring systems without understanding where current approaches break, and the layer selection bias we identify affects any practitioner deploying linear probes for safety classification.

We emphasize that our findings are diagnostic rather than prescriptive: they reveal limitations of a common evaluation methodology, not a recipe for bypassing safety measures. We encourage the community to develop layer-selection-robust probing methods before deploying hidden-state monitors in production safety systems.
